\documentclass[a4paper]{article}
\usepackage{amsfonts}
\usepackage{amsmath}
\usepackage{amssymb}
\usepackage{graphicx}     

\begin{document}
  
\noindent \large Auteur: Anna Voronovska \\
\noindent \large Studierichting: Informatica\\
\noindent \large Oefeningenreeks 7A nr. 15.1\\

\medskip

\normalsize

$f(x,y) = x-3y+5$\\

Raakvlak horizontaal als:\\

$\left\{ \begin{array}{l} \dfrac{\partial f}{\partial x} = 0 \\ \dfrac{\partial f}{\partial y} = 0 \end{array} \right.$\\

$\left\{ \begin{array}{l} \dfrac{\partial f}{\partial x} = 1 \neq 0 \\ \dfrac{\partial f}{\partial y} = -3 \neq 0 \end{array} \right.$\\

Dit stelsel heeft geen oplossing, dus er bestaat geen punt waar het raakvlak horizontaal is.

\begin{thebibliography}{999}
%\bibitem{a} Referentie
\end{thebibliography}
\end{document} 
